\section{Related Work}
\label{sec:related_work}

Our work is situated at the intersection of model merging, test-time adaptation, and optimization regularization. We review the state of the art in each of these sub-fields below.

\subsection{Weight-Space Model Merging}
Model merging has emerged as a crucial training-free method to combine multiple specialized expert models into a single multi-task system without the prohibitive computational costs of joint retraining \cite{wortsman2022model, choshen2022fusing}. A seminal work in this domain is Task Arithmetic \cite{ilharco2022editing}, which defines task vectors $\mathbf{\Delta}_k = \Theta_k - \Theta_{\text{base}}$ representing the parameter updates learned during task-specific fine-tuning. The merged model is then constructed by adding a weighted sum of these task vectors to the pre-trained base model weights: $\Theta_{\text{merged}} = \Theta_{\text{base}} + \sum_k \lambda_k \mathbf{\Delta}_k$. 

While Task Arithmetic uses a single uniform scaling factor $\lambda_k$ across all layers, subsequent works have argued that different layers of deep neural networks exhibit highly diverse sensitivities to merging and require specialized treatment. TIES-Merging \cite{yadav2023ties} resolves sign and magnitude conflicts among overlapping task vectors by keeping only top-valued parameters and aligning signs. RegMean \cite{jin2022regmean} uses out-of-distribution activation covariances computed from training samples to solve a closed-form linear projection that minimizes output representation discrepancy. Fisher-weighted merging \cite{matena2022merging} uses the diagonal of the Fisher Information Matrix to weight the contributions of different models, prioritizing parameters that are highly sensitive for specific tasks. 

\subsection{Adaptive Model Merging and Test-Time Adaptation}
Rather than relying on static, hand-crafted heuristics or requiring access to training-set statistics, recent methods have proposed \emph{adaptive model merging} that optimizes merging coefficients dynamically at test-time \cite{yang2024adamerging, jung2026symerge}. Under this setting, the merged model adapts to unlabeled test-time data streams via a self-supervised surrogate objective. 

This approach inherits foundational methodologies from Test-Time Adaptation (TTA) \cite{liang2023comprehensive, wang2020tent}. For instance, Tent \cite{wang2020tent} adapts pre-trained networks to distribution shifts by minimizing the Shannon entropy of predicted logits over unlabeled test batches, updating only the normalization layers (e.g., batch or layer normalization parameters). AdaMerging \cite{yang2024adamerging} adapts the concept of entropy minimization to weight space, directly optimizing layer-wise merging coefficients $\lambda_{k, l}$ at test-time. SyMerge \cite{jung2026symerge} extends this by identifying and leveraging cross-task synergies through a low-rank adapter-based parameterization of the merging coefficients.

While our study focuses on weight-level merging coefficient regularization within a continuous, parameter-only merging simulator, concurrent physical-only baselines like SyMerge \cite{jung2026symerge} construct physical adapter modules (e.g., low-rank adapters) that are physically injected into the network architecture during inference. Because adapter-based methods modify the forward-pass computations and require physical forward/backward operations on GPU architectures, they cannot be evaluated solely as continuous parameter-weight-space profiles within a continuous weight-merging simulator. We thus position SyMerge as an orthogonal, concurrent physical baseline that falls outside the scope of our continuous simulator, and focus our evaluation on weight-level and coefficient-level regularizers (such as TV, L2, and PolyMerge).

\subsection{Overfitting and Instability in Test-Time Adaptation}
Although test-time adaptation is highly flexible, it is notorious for optimization instability and catastrophic failure modes under realistic constraints \cite{zhao2023test}. Because TTA relies on optimizing model parameters on a single unlabeled stream, it is extremely prone to \emph{transductive overfitting}, wherein the model fits the local distribution of the adaptation batch at the expense of its global generalization capability. This issue is particularly acute when the adaptation stream is small, noisy, or contains heavy label imbalances, leading to representation collapse or catastrophic forgetting of other task capabilities \cite{yuan2023robust, niu2022towards}.

Our research reveals that this transductive overfitting is a critical, yet previously unrecognized, bottleneck in adaptive model merging. We show that unconstrained optimization of layer-wise (or projection-wise) merging coefficients gives the optimizer too many high-frequency degrees of freedom, enabling it to fit transductive noise in the test stream. While prior work has attempted to resolve TTA instability via heuristic learning rate decays or early stopping, we show that \textbf{PolyMerge} represents a more fundamental and robust solution: projecting the parameter search space onto a low-degree polynomial subspace of normalized layer depth. This physically-grounded regularization eliminates the high-frequency degrees of freedom that the optimizer exploits to overfit, providing a powerful, continuous regularizer across the depth of the network.
